\chapter{实验与结果分析}
\label{chap:experiments}

\section{概述}
本章报告课题配置 \texttt{\ThesisProfileId} 中声明的全部实验结果。下文所有数值均为由签名结果包
自动生成的宏，表格与图形亦由同一流程产出。

\input{generated/auto_tables.tex}

\ifAutoHasComparison
\section{与基线方法的对比}
表~\ref{tab:auto-comparison}~汇总了各项指标。在\MetricPrimaryName{}上，本文方法
（\MetricPrimaryProposedName）达到 \MetricPrimaryProposedMean{} $\pm$ \MetricPrimaryProposedStd{}
\MetricPrimaryUnit（置信区间 \MetricPrimaryProposedCiLow{} 至 \MetricPrimaryProposedCiHigh），
最强基线（\MetricPrimaryRunnerUpName）为 \MetricPrimaryRunnerUpMean{} $\pm$ \MetricPrimaryRunnerUpStd。
两者相差 \DiffPrimaryAbs（相对提升 \DiffPrimaryRelPct\,\%），经\StatPrimaryTest{}检验
$p$ \StatPrimaryPval（效应量 \StatPrimaryEffectName{} 为 \StatPrimaryEffect；每种方法 \StatPrimaryN{} 次运行）。
全部方法间的\StatPrimaryOmnibusTest{}给出 $p$ \StatPrimaryOmnibusPval。
\fi

\ifAutoHasAblation
\section{消融实验}
表~\ref{tab:auto-ablation}~列出了各消融配置。完整系统在\MetricPrimaryName{}上达到
\MetricPrimaryProposedMean{} $\pm$ \MetricPrimaryProposedStd{} \MetricPrimaryUnit，最优的精简配置
（\MetricPrimaryRunnerUpName）为 \MetricPrimaryRunnerUpMean，提升 \DiffPrimaryAbs（\DiffPrimaryRelPct\,\%；
\StatPrimaryTest，$p$ \StatPrimaryPval）。各配置间的\StatPrimaryOmnibusTest{}给出 $p$ \StatPrimaryOmnibusPval。
\fi

\ifAutoHasRsm
\section{响应面寻优}
各设计试验点列于表~\ref{tab:auto-design}。拟合模型（表~\ref{tab:auto-coef}）为
\begin{equation}
  \StatRsmPrimaryEquation
\end{equation}
其中 \StatRsmPrimaryFactorLegend。表~\ref{tab:auto-anova}~的方差分析给出
$R^{2}$ = \StatRsmPrimaryRsq，调整 $R^{2}$ = \StatRsmPrimaryRsqAdj（$F$ = \StatRsmPrimaryFvalue，
$p$ \StatRsmPrimaryPval），失拟检验 $p$ \StatRsmPrimaryLackOfFitPval。驻点为\StatRsmPrimaryOptNature，
位于 \StatRsmPrimaryOptPoint，预测响应为 \StatRsmPrimaryOptResponse{} \MetricPrimaryUnit。
图~\ref{fig:auto-rsm}~展示了响应面形态。
\fi

\section{图示}
\input{generated/auto_figures.tex}
